\section{Analysis}
\label{sec:analysis}

\subsection{Failure Mechanism}
\label{sec:mechanism}

Mechanical stages locate the first failing gate. They do not name a cause.
We therefore close-read artifact-level failures: the public contract, the
first-failure evaluator log, and the submitted package, plus a trajectory
screen for whether the agent inspected \texttt{repo/}. Labels follow the
protocol in Table~\ref{tab:symptoms}. Hidden test names, inputs, and
assertions are not reported. Coding is an assistant first pass (L1), not
independent human dual review, and is not a gold inter-rater study.

The census is a close-read of Pro and Flash artifact-level failures
($n=63$; Pro 28, Flash 35). We do not merge other backends into this
denominator.

\begin{table*}[t]
  \caption{Semantic symptoms used for artifact and trajectory analysis.}
  \label{tab:symptoms}
  \centering
  \begin{tabularx}{\textwidth}{@{}p{0.23\textwidth}Xp{0.27\textwidth}@{}}
    \toprule
    \textbf{Symptom} & \textbf{Observable manifestation} & \textbf{Required evidence} \\
    \midrule
    Localization & No inspection of the relevant source region
      & Trajectory + package \\
    Contract/API completion & Required symbol, member, or branch is absent
      & Public contract + package \\
    Dependency / resource closure & Helper, registry, config, or resource omitted
      & Source evidence + failing case \\
    Behavioral drift & API present; return, exception, or edge semantics differ
      & Contract + evaluator + package \\
    Packaging / modularization & Logic exists but is not independently installable
      & Package tree + isolation evidence \\
    Copy-heavy pass & Gates pass through broad vendoring
      & Passing result + RRES/copy evidence \\
    \bottomrule
  \end{tabularx}
\end{table*}

\begin{table}[t]
  \caption{Primary semantic cause on Pro and Flash artifact failures
  (L1 close-read, $n=63$). Not independent human gold.}
  \label{tab:mechanism}
  \centering
  \begin{tabular}{lrrr}
\toprule
\textbf{Primary cause} & \textbf{Total} & \textbf{Pro} & \textbf{Flash} \\
\midrule
Behavior drift & 54 & 26 & 28 \\
Contract/API completion & 7 & 2 & 5 \\
Packaging / modularization & 2 & 0 & 2 \\
Dependency / resource closure & 0 & 0 & 0 \\
Localization & 0 & 0 & 0 \\
Unknown & 0 & 0 & 0 \\
\midrule
Coded artifact failures & 63 & 28 & 35 \\
\bottomrule
\end{tabular}

\end{table}

Closure classes (API completion plus behavior drift) account for 61/63 rows.
The modal label is behavior drift: the required surface is present, but
observable semantics are not. Localization is 0 after the trajectory screen:
every remaining failure issued a deep read of \texttt{repo/} (find, grep,
cat, or file view). Packaging appears twice on Flash (an isolation
\texttt{forbidden\_imports} hit, and a submission that still imports the
upstream package after the isolation gate). We do not treat localization as
solved suite-wide; the claim is that, on this Pro+Flash artifact-fail slice,
the dominant residue is incomplete contract recovery rather than missing
packages or uninspected source trees.

A stratified L1 sample of Luna, Qwen, and OSS artifact failures
($n=45$ draws, 32 coded rows; seed
\texttt{python150-postsample-v1}) is reported in
Table~\ref{tab:postsample} and is \emph{not} pooled into the 63.
Luna remains in the same direction (drift and missing exports). Qwen and OSS
additionally fail Build by omitting transitive dependencies or by importing
the original package. That is a weaker-backend packaging story, not a
five-model pie chart.

\begin{table}[t]
  \caption{Stratified L1 sample of Luna/Qwen/OSS artifact failures
  (32 coded rows). Not merged with Table~\ref{tab:mechanism}.}
  \label{tab:postsample}
  \centering
  \begin{tabular}{lrrrr}
\toprule
\textbf{Primary cause} & \textbf{Total} & \textbf{Luna} & \textbf{Qwen} & \textbf{OSS} \\
\midrule
Behavior drift & 14 & 4 & 7 & 3 \\
Contract/API completion & 9 & 4 & 0 & 5 \\
Dependency / resource closure & 4 & 0 & 3 & 1 \\
Packaging / modularization & 5 & 0 & 1 & 4 \\
Localization & 0 & 0 & 0 & 0 \\
\midrule
Coded sample rows & 32 & 8 & 11 & 13 \\
\bottomrule
\end{tabular}

\end{table}

\paragraph{Finding 3.}
On the Pro+Flash artifact-fail slice, failures are dominated by incomplete
recovery of the required behavioral contract (a contract-closure gap), rather
than by inability to emit an installable package or by failing to inspect the
source tree. A stratified sample of the other backends does not reverse that
direction on Luna, but adds packaging and missing-helper failures on the
weaker backends. The proportions are an L1 first pass, not gold.

\subsection{Task Difficulty}
\label{sec:difficulty}

Table~\ref{tab:construction} splits Python-150 by construction cohort.
For Pro, Core-100 is 94/100 and hard3 is 21/50. Flash is 90/100 versus
18/50. Of 28 tasks unsolved by every model, 22 are hard3. OSS barely moves
(25\% versus 22\%). Entanglement level is uniformly high on this split, so
it cannot be tested. Lift type is directional but weaker
(Table~\ref{tab:lift}): Pro Direct 51/56, Adapted 55/76, Composite 9/18
(the last cell is small). Qwen's Adapted cell includes 16 empty submissions,
which are process events rather than a lift-type effect.

\begin{table}[t]
  \caption{Functional Pass by in-suite construction split on Python-150.
  hard3 is not the official Hard-50 expansion.}
  \label{tab:construction}
  \centering
  \begin{tabular}{lrr}
\toprule
\textbf{Model} & \textbf{Core-100} & \textbf{hard3} \\
\midrule
DeepSeek V4 Pro & 94/100 (94.0\%) & 21/50 (42.0\%) \\
DeepSeek V4 Flash & 90/100 (90.0\%) & 18/50 (36.0\%) \\
gpt-5.6-luna (OpenLux) & 82/100 (82.0\%) & 20/50 (40.0\%) \\
Qwen3.6-35B & 55/100 (55.0\%) & 8/50 (16.0\%) \\
GPT-OSS 120B & 25/100 (25.0\%) & 11/50 (22.0\%) \\
\bottomrule
\end{tabular}

\end{table}

\begin{table}[t]
  \caption{Functional Pass by lift type on Python-150. Composite $n=18$.}
  \label{tab:lift}
  \centering
  \begin{tabular}{lrrr}
\toprule
\textbf{Model} & \textbf{Direct} & \textbf{Adapted} & \textbf{Composite} \\
\midrule
DeepSeek V4 Pro & 51/56 (91.1\%) & 55/76 (72.4\%) & 9/18 (50.0\%) \\
DeepSeek V4 Flash & 50/56 (89.3\%) & 51/76 (67.1\%) & 7/18 (38.9\%) \\
gpt-5.6-luna (OpenLux) & 45/56 (80.4\%) & 50/76 (65.8\%) & 7/18 (38.9\%) \\
Qwen3.6-35B & 31/56 (55.4\%) & 27/76 (35.5\%) & 5/18 (27.8\%) \\
GPT-OSS 120B & 16/56 (28.6\%) & 16/76 (21.1\%) & 4/18 (22.2\%) \\
\bottomrule
\end{tabular}

\end{table}

\paragraph{Finding 4.}
Feature-lifting difficulty on Python-150 varies with the in-suite hard3
construction split and, more weakly, with lift type. Entanglement level is
uniformly high and cannot be tested here.

\subsection{Extraction Quality}
\label{sec:compactness}

Compactness is reported only on functional passes
(Table~\ref{tab:compactness}). Survivor sets differ by
model, so medians are not a second ranking of the same 150 tasks. Cross-model
comparisons use paired intersections.

\begin{table*}[t]
  \caption{Pass-conditioned compactness on freeze~v2 Python-150.
  RRES is submission LOC / reference LOC. Copy-heavy and compact are
  exclusive classes on the passing subset.}
  \label{tab:compactness}
  \centering
  \begin{tabular}{lrrrrr}
\toprule
\textbf{Model} & \textbf{$n$} & \textbf{Median RRES}
  & \textbf{Median copy} & \textbf{Copy-heavy} & \textbf{Compact} \\
\midrule
DeepSeek V4 Pro & 115 & 0.993 & 0.957 & 108 & 4 \\
DeepSeek V4 Flash & 108 & 0.998 & 0.968 & 105 & 1 \\
gpt-5.6-luna (OpenLux) & 102 & 0.815 & 0.164 & 54 & 37 \\
Qwen3.6-35B & 63 & 0.975 & 0.757 & 43 & 9 \\
GPT-OSS 120B & 36 & 0.983 & 0.180 & 18 & 12 \\
\bottomrule
\end{tabular}

\end{table*}

Pro and Flash passing packages are almost entirely copy-heavy (median copy
fraction 0.96--0.97). Luna's passing set is smaller in copy fraction (median
0.16) at median RRES 0.82. On the 97 tasks both Pro and Luna pass, median
copy is 0.97 versus 0.19. On the 24 tasks all five models pass, Pro/Flash
copy remains $\approx 0.96$--$0.97$ while Luna is 0.46. Correctness and
compactness therefore separate: a Functional Pass does not imply a compact
extraction relative to the frozen reference. Official Hard-50 has no
evaluator-side \texttt{reference\_solution/}, so we do not report RRES there.

\paragraph{Finding 5.}
Correctness and compactness are distinct: a functional pass does not imply
a compact extraction relative to the frozen reference.

\subsection{Case Studies}
\label{sec:cases}

Four cases instantiate the preceding findings. We describe observable
behavior, not Hidden test identifiers.

\paragraph{Process non-delivery.}
Qwen on \texttt{babel\_\_plural\_core\_\_001} ends without a package (tool
validation). The main-table failure is real; it is not a Hidden mismatch.

\paragraph{Public drift after inspecting the repository.}
On \texttt{alembic\_\_revision\_map\_core\_\_hard3\_001}, Pro and Flash
implement \texttt{get\_revision} but treat the argument \texttt{'base'} as
the symbolic base of the map, shadowing a revision whose identifier is
literally \texttt{base}. The API is present; the contract case is not.

\paragraph{Hidden completion with an open branch.}
On \texttt{aiohttp\_\_url\_params\_core\_\_hard3\_001}, exception types and
token checks exist, yet an invalid-name path required by the public contract
still does not raise. Localization is not the residue.

\paragraph{Isolation versus copy-heavy success.}
Flash on \texttt{typer\_\_command\_parser\_core\_\_001} fails Isolation via
\texttt{forbidden\_imports}. On \texttt{blinker\_\_signal\_registry\_core\_\_001},
Pro, Flash, and Luna all pass with copy-heavy packages (RRES 9.0 / 18.4 /
4.6). Passing is not compact.

\subsection{Information Boundary and Cost (Non-Main)}
\label{sec:nonmain}

Two older measurements diagnose mechanism and are not freeze~v2 Main scores.

\paragraph{Public feedback.}
A paired ablation on 12 DeepSeek V4 Flash tasks (historical slice) mounts
\texttt{public\_tests/} while keeping Hidden private. Main is 0/12;
Public-feedback is 4/12 (Table~\ref{tab:public-feedback}). All six selected
Public failures flip Public to 1; Hidden usually does not. Executable Public
tests are a real bottleneck and not a general Hidden oracle. The slice is not
the Python-150 leaderboard.

\begin{table}[t]
  \caption{Paired Public-feedback ablation on 12 historical Flash tasks.
  Not freeze~v2 Main.}
  \label{tab:public-feedback}
  \centering
  \begin{tabular}{lr}
    \toprule
    \textbf{Condition} & \textbf{Functional Pass} \\
    \midrule
    Full-Repository / No-Hint Main & 0/12 \\
    Public-feedback & 4/12 \\
    \bottomrule
  \end{tabular}
\end{table}

\paragraph{Earliest sufficiency $T^*$.}
On replayable historical Flash passing trajectories (superseded 150+E50
suite), 138 of 145 passes have gold. Median $T^*/T_{\mathrm{total}}$ is 0.40,
with a median 0.75M tokens after a sufficient package first exists
(Table~\ref{tab:tstar}). Post-sufficiency spend is mostly self-testing the
agent cannot ground in Hidden. This is not a stopping rule and is not
recomputed on freeze~v2.

\begin{table}[t]
  \caption{Earliest functional sufficiency on historical DeepSeek V4 Flash
  passing trajectories (not freeze~v2).}
  \label{tab:tstar}
  \centering
  \begin{tabular}{lrrr}
    \toprule
    \textbf{Slice} & \textbf{$n$} & \textbf{Median $T^*/T$}
      & \textbf{Post-$T^*$ tokens} \\
    \midrule
    All gold passes & 138 & 0.40 & 0.75M \\
    Direct & 58 & 0.36 & 0.74M \\
    Adapted & 52 & 0.40 & 0.79M \\
    Composite & 28 & 0.51 & 0.70M \\
    \bottomrule
  \end{tabular}
\end{table}

\paragraph{Process scaffolds.}
Development interventions (self-generated probes, test-first extraction,
repair loops, checkpoint replay, 2M caps) did not produce a stable Functional
gain over legal Main on the subsets where they were tried
(Table~\ref{tab:negative}). They are negative diagnostics, not methods in
the leaderboard.

\begin{table*}[t]
  \caption{Development interventions and their diagnostic outcomes.
  Not freeze~v2 Main.}
  \label{tab:negative}
  \centering
  \begin{tabularx}{\textwidth}{@{}p{0.25\textwidth}p{0.32\textwidth}X@{}}
    \toprule
    \textbf{Hypothesis} & \textbf{Interventions} & \textbf{Observed outcome} \\
    \midrule
    Better self-verification is sufficient
      & Self-generated probes, test-first extraction, verification ledger
      & No stable Functional gain; some variants cost more or regress \\
    Executable contract checklists close behavior
      & Self-Contract, Exec-Contract, spec-adversarial checks
      & Can repair Public; evaluated Hidden subsets do not improve reliably \\
    Correct intermediate packages are overwritten
      & Repair loops, Rescue+, best-so-far replay
      & No stable gain; 0/51 failed trajectories contain an earlier pass \\
    Tighter budget/context control improves quality
      & 2M cap, adaptive stop, compressed context, structure guidance
      & Token savings trade off against Functional Pass \\
    \bottomrule
  \end{tabularx}
\end{table*}
